\documentclass[12pt,a4paper]{article}

\usepackage[T2A]{fontenc}
\usepackage[utf8]{inputenc}
\usepackage[russian]{babel}
\usepackage{amsmath,amssymb,amsthm}
\usepackage{geometry}
\usepackage{listings}
\usepackage{xcolor}
\usepackage{hyperref}
\usepackage{array}
\usepackage{enumitem}

\geometry{margin=2cm}

\lstdefinestyle{code}{
    basicstyle=\ttfamily\small,
    keywordstyle=\color{blue},
    commentstyle=\color{gray},
    stringstyle=\color{teal},
    showstringspaces=false,
    breaklines=true,
    frame=single,
    tabsize=2,
    literate=
      {А}{{\CYRA}}1
      {Б}{{\CYRB}}1
      {В}{{\CYRV}}1
      {Г}{{\CYRG}}1
      {Д}{{\CYRD}}1
      {Е}{{\CYRE}}1
      {Ё}{{\CYRYO}}1
      {Ж}{{\CYRZH}}1
      {З}{{\CYRZ}}1
      {И}{{\CYRI}}1
      {Й}{{\CYRISHRT}}1
      {К}{{\CYRK}}1
      {Л}{{\CYRL}}1
      {М}{{\CYRM}}1
      {Н}{{\CYRN}}1
      {О}{{\CYRO}}1
      {П}{{\CYRP}}1
      {Р}{{\CYRR}}1
      {С}{{\CYRS}}1
      {Т}{{\CYRT}}1
      {У}{{\CYRU}}1
      {Ф}{{\CYRF}}1
      {Х}{{\CYRH}}1
      {Ц}{{\CYRC}}1
      {Ч}{{\CYRCH}}1
      {Ш}{{\CYRSH}}1
      {Щ}{{\CYRSHCH}}1
      {Ъ}{{\CYRHRDSN}}1
      {Ы}{{\CYRERY}}1
      {Ь}{{\CYRSFTSN}}1
      {Э}{{\CYREREV}}1
      {Ю}{{\CYRYU}}1
      {Я}{{\CYRYA}}1
      {а}{{\cyra}}1
      {б}{{\cyrb}}1
      {в}{{\cyrv}}1
      {г}{{\cyrg}}1
      {д}{{\cyrd}}1
      {е}{{\cyre}}1
      {ё}{{\cyryo}}1
      {ж}{{\cyrzh}}1
      {з}{{\cyrz}}1
      {и}{{\cyri}}1
      {й}{{\cyrishrt}}1
      {к}{{\cyrk}}1
      {л}{{\cyrl}}1
      {м}{{\cyrm}}1
      {н}{{\cyrn}}1
      {о}{{\cyro}}1
      {п}{{\cyrp}}1
      {р}{{\cyrr}}1
      {с}{{\cyrs}}1
      {т}{{\cyrt}}1
      {у}{{\cyru}}1
      {ф}{{\cyrf}}1
      {х}{{\cyrh}}1
      {ц}{{\cyrc}}1
      {ч}{{\cyrch}}1
      {ш}{{\cyrsh}}1
      {щ}{{\cyrshch}}1
      {ъ}{{\cyrhrdsn}}1
      {ы}{{\cyrery}}1
      {ь}{{\cyrsftsn}}1
      {э}{{\cyrerev}}1
      {ю}{{\cyryu}}1
      {я}{{\cyrya}}1
}

\lstset{style=code}

\title{Билет №55 \\ \small Системы сборки ПО и управления модулями. Универсальные системы (Make) и ориентированные на конкретные языки (CMake, Cargo, Maven и т.д.) (СпБГУ)}
\author{Сериков Владислав Александрович}
\date{16 мая 2026}

\begin{document}

\maketitle
\tableofcontents
\newpage

\section{Общие понятия}

\subsection{Проблема сборки программного обеспечения}

Сборка программного обеспечения --- это процесс получения исполняемых файлов, библиотек, пакетов или иных артефактов из исходных текстов, ресурсов и зависимостей проекта.

Типичный процесс сборки включает:
\begin{enumerate}
    \item предварительную обработку исходных файлов;
    \item компиляцию;
    \item генерацию кода;
    \item линковку;
    \item запуск тестов;
    \item упаковку;
    \item установку;
    \item публикацию артефактов.
\end{enumerate}

Например, для программы на C процесс может иметь вид:
\[
\texttt{main.c} \longrightarrow \texttt{main.o} \longrightarrow \texttt{app}
\]

Здесь:
\begin{itemize}
    \item \texttt{main.c} --- исходный файл;
    \item \texttt{main.o} --- объектный файл;
    \item \texttt{app} --- исполняемый файл.
\end{itemize}

Если проект состоит из большого числа файлов, ручной запуск команд компиляции становится неудобным и ненадёжным. Поэтому используются системы сборки.

\subsection{Система сборки}

\textbf{Система сборки} --- это программный инструмент, который автоматизирует процесс преобразования исходных файлов проекта в целевые артефакты.

Основные задачи системы сборки:
\begin{enumerate}
    \item описание зависимостей между файлами и целями;
    \item определение порядка выполнения команд;
    \item инкрементальная пересборка только изменившихся частей проекта;
    \item поддержка разных конфигураций сборки;
    \item интеграция с компиляторами, линковщиками, тестовыми фреймворками;
    \item генерация и установка артефактов.
\end{enumerate}

\subsection{Модуль и зависимость}

\textbf{Модуль} --- логически обособленная часть программы, которая имеет собственный интерфейс и может быть использована другими частями программы.

В разных языках модулем могут быть:
\begin{itemize}
    \item файл исходного кода;
    \item пакет;
    \item библиотека;
    \item crate в Rust;
    \item artifact в Maven;
    \item npm-пакет в JavaScript;
    \item Python-пакет.
\end{itemize}

\textbf{Зависимость} --- отношение, при котором один модуль, файл или артефакт требует наличия другого для корректной сборки или выполнения.

Пример зависимости:
\[
A \rightarrow B
\]
означает, что модуль $A$ зависит от модуля $B$.

\subsection{Граф зависимостей}

Многие системы сборки представляют проект в виде ориентированного графа зависимостей.

\textbf{Граф зависимостей} --- это ориентированный граф
\[
G = (V, E),
\]
где:
\begin{itemize}
    \item $V$ --- множество целей, файлов или модулей;
    \item $E$ --- множество ориентированных рёбер зависимости.
\end{itemize}

Если существует ребро
\[
u \rightarrow v,
\]
то цель $u$ зависит от цели $v$, то есть перед сборкой $u$ необходимо убедиться, что $v$ уже собрана или актуальна.

Пример:
\[
\texttt{app} \rightarrow \texttt{main.o}, \quad
\texttt{app} \rightarrow \texttt{util.o}, \quad
\texttt{main.o} \rightarrow \texttt{main.c}
\]

Графически:

\[
\begin{array}{c}
\texttt{main.c} \quad \texttt{util.c} \\
\downarrow \quad \downarrow \\
\texttt{main.o} \quad \texttt{util.o} \\
\quad \searrow \quad \swarrow \\
\texttt{app}
\end{array}
\]

\section{Основные свойства систем сборки}

\subsection{Инкрементальность}

\textbf{Инкрементальная сборка} --- это сборка, при которой повторно выполняются только те действия, результаты которых могли устареть после изменения входных данных.

Например, если изменился только файл \texttt{main.c}, то нет необходимости заново компилировать \texttt{util.c}. Достаточно пересобрать:
\[
\texttt{main.c} \rightarrow \texttt{main.o} \rightarrow \texttt{app}
\]

\subsection{Воспроизводимость}

\textbf{Воспроизводимая сборка} --- это такая сборка, которая при одинаковых исходных данных, зависимостях, параметрах и окружении даёт одинаковый результат.

Воспроизводимость важна для:
\begin{itemize}
    \item отладки;
    \item безопасности;
    \item непрерывной интеграции;
    \item распространения пакетов;
    \item проверки, что бинарный артефакт действительно соответствует исходному коду.
\end{itemize}

Факторы, нарушающие воспроизводимость:
\begin{enumerate}
    \item незафиксированные версии зависимостей;
    \item использование текущего времени в артефактах;
    \item зависимость от локальных путей;
    \item недетерминированный порядок обхода файлов;
    \item различия в версиях компиляторов;
    \item сетевые обращения во время сборки.
\end{enumerate}

\subsection{Параллельность}

Если две цели не зависят друг от друга, их можно собирать параллельно.

Пусть есть зависимости:
\[
\texttt{app} \rightarrow \texttt{a.o}, \quad
\texttt{app} \rightarrow \texttt{b.o}
\]

Если \texttt{a.o} и \texttt{b.o} не зависят друг от друга, то они могут быть собраны одновременно.

Это позволяет ускорить сборку на многоядерных процессорах.

\subsection{Кэширование}

\textbf{Кэширование сборки} --- это сохранение результатов ранее выполненных действий для повторного использования.

Идея:
\[
\text{ключ} = H(\text{исходные файлы}, \text{команды}, \text{параметры}, \text{зависимости})
\]

Если ключ уже есть в кэше, действие можно не выполнять заново.

\section{Алгоритмическая основа систем сборки}

\subsection{Топологическая сортировка}

Для определения порядка сборки используется топологическая сортировка графа зависимостей.

\textbf{Топологическая сортировка} ориентированного графа --- это такой линейный порядок его вершин, что для каждого ребра
\[
u \rightarrow v
\]
вершина $v$ располагается раньше вершины $u$, если $u$ зависит от $v$.

\subsubsection{Алгоритм топологической сортировки на основе входящих степеней}

Пусть дан ориентированный ациклический граф зависимостей $G=(V,E)$.

\textbf{Шаги алгоритма:}
\begin{enumerate}
    \item Для каждой вершины вычислить количество входящих рёбер.
    \item Поместить в очередь все вершины с нулевой входящей степенью.
    \item Пока очередь не пуста:
    \begin{enumerate}
        \item извлечь вершину $v$;
        \item добавить $v$ в результат;
        \item для каждого ребра $v \rightarrow u$ уменьшить входящую степень $u$ на единицу;
        \item если входящая степень $u$ стала нулевой, добавить $u$ в очередь.
    \end{enumerate}
    \item Если в результате оказалось меньше вершин, чем в графе, значит, в графе есть цикл.
\end{enumerate}

\subsubsection{Пример}

Пусть:
\[
A \rightarrow B, \quad A \rightarrow C, \quad B \rightarrow D, \quad C \rightarrow D.
\]

Это означает, что $A$ зависит от $B$ и $C$, а $B$ и $C$ зависят от $D$.

Правильный порядок сборки:
\[
D, B, C, A
\]
или
\[
D, C, B, A.
\]

\subsubsection{Минимальная реализация на псевдокоде}

\begin{lstlisting}
topological_sort(graph):
    indeg = map from vertex to 0

    for v in graph.vertices:
        for u in graph[v]:
            indeg[u] += 1

    queue = all vertices with indeg[v] == 0
    result = []

    while queue is not empty:
        v = queue.pop()
        result.append(v)

        for u in graph[v]:
            indeg[u] -= 1
            if indeg[u] == 0:
                queue.push(u)

    if len(result) != len(graph.vertices):
        error "cycle detected"

    return result
\end{lstlisting}

\subsection{Обнаружение устаревших целей}

Классические системы сборки, например Make, часто используют временные метки файлов.

Цель считается устаревшей, если:
\[
mtime(\text{target}) < \max_i mtime(\text{dependency}_i),
\]
где $mtime(x)$ --- время последнего изменения файла $x$.

Иначе говоря, если хотя бы одна зависимость новее цели, цель нужно пересобрать.

\subsubsection{Алгоритм инкрементальной пересборки по временным меткам}

Пусть дана цель $T$.

\textbf{Шаги:}
\begin{enumerate}
    \item Если $T$ не существует, пометить её как требующую сборки.
    \item Рекурсивно проверить все зависимости $D_1,\ldots,D_n$.
    \item Если хотя бы одна зависимость была пересобрана, проверить необходимость сборки $T$.
    \item Если
    \[
    mtime(T) < \max_i mtime(D_i),
    \]
    выполнить команду сборки $T$.
    \item Иначе считать $T$ актуальной.
\end{enumerate}

\subsubsection{Пример}

Пусть:
\[
mtime(\texttt{main.c}) = 10:00,
\]
\[
mtime(\texttt{main.o}) = 09:55.
\]

Так как исходный файл новее объектного:
\[
10:00 > 09:55,
\]
то \texttt{main.o} нужно пересобрать.

\subsection{Циклические зависимости}

Циклическая зависимость возникает, если:
\[
A \rightarrow B \rightarrow C \rightarrow A.
\]

В таком случае невозможно выбрать корректный порядок сборки: каждая цель требует предварительной сборки другой.

Большинство систем сборки либо запрещают циклы, либо требуют явно разрывать их с помощью специальных механизмов.

\section{Make}

\subsection{Назначение Make}

\textbf{Make} --- универсальная система сборки, основанная на правилах, целях, зависимостях и командах.

Make особенно часто используется для проектов на C и C++, но не привязан к конкретному языку программирования. С его помощью можно автоматизировать любые действия, которые выражаются командами оболочки.

\subsection{Makefile}

Описание сборки помещается в файл \texttt{Makefile}.

Общая форма правила:

\begin{lstlisting}
target: dependencies
	command
\end{lstlisting}

Важно: команда в классическом Makefile должна начинаться с символа табуляции.

\subsection{Цели}

\textbf{Цель} --- это объект, который нужно построить. Обычно целью является файл, например объектный файл или исполняемый файл.

Пример:

\begin{lstlisting}
app: main.o util.o
	gcc main.o util.o -o app
\end{lstlisting}

Здесь:
\begin{itemize}
    \item \texttt{app} --- цель;
    \item \texttt{main.o} и \texttt{util.o} --- зависимости;
    \item \texttt{gcc main.o util.o -o app} --- команда сборки.
\end{itemize}

\subsection{Минимальный пример Makefile}

Пусть есть файлы:
\begin{itemize}
    \item \texttt{main.c};
    \item \texttt{util.c};
    \item \texttt{util.h}.
\end{itemize}

\begin{lstlisting}
CC = gcc
CFLAGS = -Wall -O2

app: main.o util.o
	$(CC) main.o util.o -o app

main.o: main.c util.h
	$(CC) $(CFLAGS) -c main.c

util.o: util.c util.h
	$(CC) $(CFLAGS) -c util.c

clean:
	rm -f *.o app
\end{lstlisting}

Сборка:
\begin{lstlisting}
make
\end{lstlisting}

Очистка:
\begin{lstlisting}
make clean
\end{lstlisting}

\subsection{Фиктивные цели}

Некоторые цели не соответствуют файлам. Например, \texttt{clean} обычно не создаёт файл с именем \texttt{clean}. Такие цели называют фиктивными.

Их объявляют с помощью \texttt{.PHONY}:

\begin{lstlisting}
.PHONY: clean

clean:
	rm -f *.o app
\end{lstlisting}

Если не указать \texttt{.PHONY}, то наличие файла \texttt{clean} в каталоге может помешать выполнению команды.

\subsection{Переменные в Make}

Make поддерживает переменные:

\begin{lstlisting}
CC = gcc
CFLAGS = -Wall -O2
\end{lstlisting}

Использование:

\begin{lstlisting}
$(CC) $(CFLAGS) -c main.c
\end{lstlisting}

Это позволяет централизованно менять параметры сборки.

\subsection{Автоматические переменные}

В Make есть автоматические переменные:

\begin{center}
\begin{tabular}{|c|p{9cm}|}
\hline
Переменная & Значение \\
\hline
\texttt{\$@} & имя цели \\
\hline
\texttt{\$<} & первая зависимость \\
\hline
\texttt{\$\string^} & все зависимости \\
\hline
\end{tabular}
\end{center}

Пример:

\begin{lstlisting}
app: main.o util.o
	gcc $^ -o $@
\end{lstlisting}

Эквивалентно:

\begin{lstlisting}
gcc main.o util.o -o app
\end{lstlisting}

\subsection{Шаблонные правила}

Шаблонное правило:

\begin{lstlisting}
%.o: %.c
	gcc -Wall -O2 -c $< -o $@
\end{lstlisting}

Это означает, что любой файл \texttt{name.o} может быть построен из \texttt{name.c}.

Полный пример:

\begin{lstlisting}
CC = gcc
CFLAGS = -Wall -O2
OBJS = main.o util.o

app: $(OBJS)
	$(CC) $(OBJS) -o app

%.o: %.c
	$(CC) $(CFLAGS) -c $< -o $@

.PHONY: clean
clean:
	rm -f $(OBJS) app
\end{lstlisting}

\subsection{Параллельная сборка}

Make поддерживает параллельное выполнение независимых правил:

\begin{lstlisting}
make -j4
\end{lstlisting}

Здесь \texttt{-j4} означает использование до четырёх параллельных задач.

\subsection{Преимущества Make}

\begin{enumerate}
    \item Универсальность.
    \item Простая модель: цели, зависимости, команды.
    \item Широкая распространённость.
    \item Хорошая поддержка инкрементальной сборки.
    \item Возможность использовать любые команды оболочки.
\end{enumerate}

\subsection{Недостатки Make}

\begin{enumerate}
    \item Слабая переносимость между платформами.
    \item Сложность сопровождения больших Makefile.
    \item Зависимость от командной оболочки.
    \item Ограниченная встроенная поддержка модулей и внешних библиотек.
    \item Неудобство описания сложных конфигураций.
\end{enumerate}

\section{CMake}

\subsection{Назначение CMake}

\textbf{CMake} --- это кроссплатформенная система конфигурации и генерации сборочных файлов. CMake не всегда непосредственно выполняет сборку, а часто генерирует файлы для других систем: Make, Ninja, Visual Studio, Xcode и других.

Типичный процесс:

\[
\texttt{CMakeLists.txt}
\longrightarrow
\text{файлы конкретной системы сборки}
\longrightarrow
\text{артефакты}
\]

Например:

\[
\texttt{CMakeLists.txt}
\longrightarrow
\texttt{Makefile}
\longrightarrow
\texttt{app}
\]

\subsection{Основные понятия CMake}

Основными сущностями CMake являются:
\begin{itemize}
    \item проект;
    \item цель;
    \item исполняемый файл;
    \item библиотека;
    \item свойства цели;
    \item зависимости между целями;
    \item конфигурация сборки.
\end{itemize}

\subsection{Минимальный пример CMake}

Файл \texttt{CMakeLists.txt}:

\begin{lstlisting}
cmake_minimum_required(VERSION 3.20)

project(Example C)

add_executable(app main.c util.c)
\end{lstlisting}

Сборка:

\begin{lstlisting}
cmake -S . -B build
cmake --build build
\end{lstlisting}

Здесь:
\begin{itemize}
    \item \texttt{-S .} указывает каталог с исходниками;
    \item \texttt{-B build} указывает каталог сборки;
    \item \texttt{cmake --build build} запускает сборку сгенерированной системой.
\end{itemize}

\subsection{Out-of-source build}

Хорошей практикой является сборка вне каталога исходного кода:

\begin{lstlisting}
project/
  CMakeLists.txt
  main.c
  util.c
  build/
\end{lstlisting}

Преимущества:
\begin{enumerate}
    \item исходный каталог не засоряется временными файлами;
    \item можно иметь несколько конфигураций сборки;
    \item проще удалить результаты сборки.
\end{enumerate}

\subsection{Цели в CMake}

CMake ориентирован на цели.

Пример библиотеки и исполняемого файла:

\begin{lstlisting}
cmake_minimum_required(VERSION 3.20)

project(Example C)

add_library(util util.c)
add_executable(app main.c)

target_link_libraries(app PRIVATE util)
\end{lstlisting}

Здесь:
\begin{itemize}
    \item \texttt{util} --- библиотека;
    \item \texttt{app} --- исполняемый файл;
    \item \texttt{app} зависит от \texttt{util}.
\end{itemize}

\subsection{Области видимости зависимостей}

Команда \texttt{target\_link\_libraries} использует ключевые слова:
\begin{itemize}
    \item \texttt{PRIVATE};
    \item \texttt{PUBLIC};
    \item \texttt{INTERFACE}.
\end{itemize}

Их смысл:

\begin{center}
\begin{tabular}{|c|p{10cm}|}
\hline
Ключевое слово & Смысл \\
\hline
\texttt{PRIVATE} & Зависимость нужна только для сборки данной цели \\
\hline
\texttt{PUBLIC} & Зависимость нужна данной цели и её пользователям \\
\hline
\texttt{INTERFACE} & Зависимость нужна только пользователям цели \\
\hline
\end{tabular}
\end{center}

Пример:

\begin{lstlisting}
add_library(mathlib mathlib.c)
target_include_directories(mathlib PUBLIC include)

add_executable(app main.c)
target_link_libraries(app PRIVATE mathlib)
\end{lstlisting}

Если \texttt{mathlib} публикует заголовочные файлы, то они становятся доступны целям, которые с ней связываются.

\subsection{Поиск внешних библиотек}

CMake предоставляет команду \texttt{find\_package}:

\begin{lstlisting}
find_package(OpenSSL REQUIRED)

add_executable(app main.c)
target_link_libraries(app PRIVATE OpenSSL::SSL OpenSSL::Crypto)
\end{lstlisting}

Современный CMake предпочитает импортированные цели вида:
\[
\texttt{Package::Target}.
\]

\subsection{Преимущества CMake}

\begin{enumerate}
    \item Кроссплатформенность.
    \item Поддержка разных генераторов.
    \item Удобная модель целей.
    \item Хорошая интеграция с C и C++.
    \item Поддержка тестирования, установки, упаковки.
    \item Возможность экспорта библиотек для использования другими проектами.
\end{enumerate}

\subsection{Недостатки CMake}

\begin{enumerate}
    \item Сложный язык конфигурации.
    \item Исторически накопившееся множество стилей написания.
    \item Не является полноценным менеджером пакетов.
    \item Ошибки конфигурации могут проявляться только на этапе генерации или сборки.
\end{enumerate}

\section{Ninja}

\subsection{Назначение Ninja}

\textbf{Ninja} --- минималистичная система сборки, ориентированная на скорость.

В отличие от Make, Ninja обычно не пишут вручную. Чаще его файлы генерируются более высокоуровневыми инструментами, например CMake или Meson.

\subsection{Пример файла build.ninja}

\begin{lstlisting}
rule cc
  command = gcc -c $in -o $out

rule link
  command = gcc $in -o $out

build main.o: cc main.c
build util.o: cc util.c
build app: link main.o util.o
\end{lstlisting}

Сборка:

\begin{lstlisting}
ninja
\end{lstlisting}

\subsection{Преимущества Ninja}

\begin{enumerate}
    \item высокая скорость запуска;
    \item эффективная инкрементальная сборка;
    \item простая внутренняя модель;
    \item удобство как backend для других систем.
\end{enumerate}

\subsection{Недостатки Ninja}

\begin{enumerate}
    \item неудобен для ручного написания сложных проектов;
    \item не является менеджером пакетов;
    \item практически не решает задачи конфигурации проекта.
\end{enumerate}

\section{Cargo}

\subsection{Назначение Cargo}

\textbf{Cargo} --- стандартная система сборки и менеджер пакетов для языка Rust.

Cargo выполняет несколько ролей:
\begin{enumerate}
    \item создаёт структуру проекта;
    \item собирает проект;
    \item запускает тесты;
    \item управляет зависимостями;
    \item генерирует документацию;
    \item публикует пакеты;
    \item поддерживает lock-файл для воспроизводимости.
\end{enumerate}

\subsection{Crate и package}

В экосистеме Rust важны два понятия:
\begin{itemize}
    \item \textbf{crate} --- единица компиляции Rust;
    \item \textbf{package} --- пакет, содержащий один или несколько crate.
\end{itemize}

Пакет описывается файлом:
\[
\texttt{Cargo.toml}.
\]

\subsection{Создание проекта}

\begin{lstlisting}
cargo new hello
cd hello
cargo run
\end{lstlisting}

Структура:

\begin{lstlisting}
hello/
  Cargo.toml
  src/
    main.rs
\end{lstlisting}

\subsection{Пример Cargo.toml}

\begin{lstlisting}
[package]
name = "hello"
version = "0.1.0"
edition = "2021"

[dependencies]
serde = "1.0"
\end{lstlisting}

\subsection{Основные команды Cargo}

\begin{center}
\begin{tabular}{|c|p{9cm}|}
\hline
Команда & Назначение \\
\hline
\texttt{cargo build} & собрать проект \\
\hline
\texttt{cargo run} & собрать и запустить \\
\hline
\texttt{cargo test} & запустить тесты \\
\hline
\texttt{cargo check} & проверить проект без полной генерации бинарного кода \\
\hline
\texttt{cargo doc} & сгенерировать документацию \\
\hline
\texttt{cargo publish} & опубликовать пакет \\
\hline
\end{tabular}
\end{center}

\subsection{Cargo.lock}

Файл \texttt{Cargo.lock} фиксирует точные версии зависимостей, выбранные решателем зависимостей.

Назначение:
\begin{itemize}
    \item обеспечить воспроизводимость сборки;
    \item зафиксировать транзитивные зависимости;
    \item предотвратить неожиданные обновления.
\end{itemize}

Для исполняемых приложений \texttt{Cargo.lock} обычно коммитят в систему контроля версий. Для библиотек подход зависит от политики проекта.

\subsection{Преимущества Cargo}

\begin{enumerate}
    \item тесная интеграция с языком Rust;
    \item единый стандарт структуры проекта;
    \item встроенный менеджер зависимостей;
    \item удобная работа с тестами и документацией;
    \item lock-файл для воспроизводимости;
    \item простые команды.
\end{enumerate}

\subsection{Недостатки Cargo}

\begin{enumerate}
    \item ориентирован почти исключительно на Rust;
    \item сложные сценарии интеграции с C/C++ требуют дополнительных build script;
    \item не заменяет системный пакетный менеджер.
\end{enumerate}

\section{Maven}

\subsection{Назначение Maven}

\textbf{Apache Maven} --- система сборки и управления зависимостями, прежде всего для проектов на Java и JVM-языках.

Maven использует декларативный подход: разработчик описывает проект в файле \texttt{pom.xml}, а Maven выполняет сборку согласно стандартному жизненному циклу.

\subsection{POM}

\textbf{POM} --- Project Object Model, модель проекта Maven.

Файл \texttt{pom.xml} содержит:
\begin{itemize}
    \item координаты проекта;
    \item зависимости;
    \item плагины;
    \item настройки сборки;
    \item информацию о репозиториях;
    \item метаданные проекта.
\end{itemize}

Минимальный пример:

\begin{lstlisting}[language=XML]
<project xmlns="http://maven.apache.org/POM/4.0.0"
         xmlns:xsi="http://www.w3.org/2001/XMLSchema-instance"
         xsi:schemaLocation="http://maven.apache.org/POM/4.0.0
                             https://maven.apache.org/xsd/maven-4.0.0.xsd">

    <modelVersion>4.0.0</modelVersion>

    <groupId>org.example</groupId>
    <artifactId>demo</artifactId>
    <version>1.0.0</version>

</project>
\end{lstlisting}

\subsection{Координаты артефакта}

Артефакт Maven идентифицируется координатами:
\[
(\texttt{groupId}, \texttt{artifactId}, \texttt{version})
\]

Например:

\begin{lstlisting}[language=XML]
<groupId>org.apache.commons</groupId>
<artifactId>commons-lang3</artifactId>
<version>3.14.0</version>
\end{lstlisting}

\subsection{Стандартная структура проекта Maven}

\begin{lstlisting}
project/
  pom.xml
  src/
    main/
      java/
      resources/
    test/
      java/
      resources/
\end{lstlisting}

Принцип Maven:
\[
\text{convention over configuration}
\]
то есть соглашения важнее явной конфигурации.

\subsection{Жизненный цикл Maven}

Maven определяет фазы жизненного цикла сборки.

Основные фазы:
\begin{center}
\begin{tabular}{|c|p{9cm}|}
\hline
Фаза & Назначение \\
\hline
\texttt{validate} & проверить корректность проекта \\
\hline
\texttt{compile} & скомпилировать исходный код \\
\hline
\texttt{test} & запустить тесты \\
\hline
\texttt{package} & упаковать результат, например в JAR \\
\hline
\texttt{verify} & выполнить дополнительные проверки \\
\hline
\texttt{install} & установить артефакт в локальный репозиторий \\
\hline
\texttt{deploy} & опубликовать артефакт в удалённый репозиторий \\
\hline
\end{tabular}
\end{center}

Команда:

\begin{lstlisting}
mvn package
\end{lstlisting}

выполняет не только фазу \texttt{package}, но и предшествующие фазы:
\[
\texttt{validate} \rightarrow \texttt{compile} \rightarrow \texttt{test} \rightarrow \texttt{package}.
\]

\subsection{Зависимости в Maven}

Пример зависимости:

\begin{lstlisting}[language=XML]
<dependencies>
    <dependency>
        <groupId>org.apache.commons</groupId>
        <artifactId>commons-lang3</artifactId>
        <version>3.14.0</version>
    </dependency>
</dependencies>
\end{lstlisting}

Maven загружает зависимости из репозиториев, например из Maven Central, и помещает их в локальный кэш.

\subsection{Области видимости зависимостей}

Maven поддерживает scopes:

\begin{center}
\begin{tabular}{|c|p{9cm}|}
\hline
Scope & Смысл \\
\hline
\texttt{compile} & зависимость нужна для компиляции и выполнения \\
\hline
\texttt{provided} & нужна для компиляции, но будет предоставлена окружением \\
\hline
\texttt{runtime} & нужна только во время выполнения \\
\hline
\texttt{test} & нужна только для тестов \\
\hline
\texttt{system} & указывается явно локальный путь к JAR \\
\hline
\texttt{import} & используется для импорта dependency management в BOM \\
\hline
\end{tabular}
\end{center}

Пример тестовой зависимости:

\begin{lstlisting}[language=XML]
<dependency>
    <groupId>org.junit.jupiter</groupId>
    <artifactId>junit-jupiter</artifactId>
    <version>5.10.0</version>
    <scope>test</scope>
</dependency>
\end{lstlisting}

\subsection{Транзитивные зависимости}

Если проект $A$ зависит от $B$, а $B$ зависит от $C$, то $A$ косвенно зависит от $C$:
\[
A \rightarrow B \rightarrow C.
\]

Maven автоматически учитывает такие транзитивные зависимости.

\subsection{Разрешение конфликтов версий}

Если одна и та же библиотека встречается в графе зависимостей в разных версиях, Maven выбирает версию по правилам разрешения конфликтов.

Одно из основных правил:
\begin{quote}
при конфликте выбирается версия, находящаяся ближе к корню графа зависимостей.
\end{quote}

Пример:

\[
A \rightarrow B \rightarrow D:2.0
\]
\[
A \rightarrow E \rightarrow F \rightarrow D:1.0
\]

Версия \texttt{D:2.0} ближе к $A$, поэтому она будет выбрана.

\subsection{Преимущества Maven}

\begin{enumerate}
    \item стандартная структура проекта;
    \item развитая модель зависимостей;
    \item большое количество плагинов;
    \item интеграция с Maven Central;
    \item поддержка жизненного цикла сборки;
    \item удобен для Java-проектов.
\end{enumerate}

\subsection{Недостатки Maven}

\begin{enumerate}
    \item XML-конфигурация может быть громоздкой;
    \item сложная модель транзитивных зависимостей;
    \item ограниченная гибкость по сравнению с более программируемыми системами;
    \item возможны неожиданные конфликты версий.
\end{enumerate}

\section{Gradle}

\subsection{Назначение Gradle}

\textbf{Gradle} --- система сборки, широко применяемая в Java, Kotlin, Android и других JVM-проектах.

Gradle сочетает:
\begin{itemize}
    \item декларативное описание зависимостей;
    \item программируемые build-скрипты;
    \item поддержку многомодульных проектов;
    \item инкрементальную сборку;
    \item кэширование;
    \item плагины.
\end{itemize}

\subsection{Пример build.gradle.kts}

\begin{lstlisting}
plugins {
    application
}

repositories {
    mavenCentral()
}

dependencies {
    implementation("org.apache.commons:commons-lang3:3.14.0")
    testImplementation("org.junit.jupiter:junit-jupiter:5.10.0")
}

application {
    mainClass.set("org.example.Main")
}
\end{lstlisting}

\subsection{Конфигурации зависимостей}

Типичные конфигурации:
\begin{center}
\begin{tabular}{|c|p{9cm}|}
\hline
Конфигурация & Смысл \\
\hline
\texttt{implementation} & зависимость нужна для реализации, но не экспортируется пользователям \\
\hline
\texttt{api} & зависимость является частью публичного API \\
\hline
\texttt{compileOnly} & нужна только при компиляции \\
\hline
\texttt{runtimeOnly} & нужна только во время выполнения \\
\hline
\texttt{testImplementation} & нужна для тестов \\
\hline
\end{tabular}
\end{center}

\subsection{Преимущества Gradle}

\begin{enumerate}
    \item высокая гибкость;
    \item удобство для многомодульных проектов;
    \item хорошая поддержка Android;
    \item инкрементальность и кэширование;
    \item Kotlin DSL и Groovy DSL.
\end{enumerate}

\subsection{Недостатки Gradle}

\begin{enumerate}
    \item сложность модели;
    \item возможная недетерминированность при неаккуратных скриптах;
    \item порог входа выше, чем у Maven;
    \item долгий этап конфигурации в больших проектах без оптимизаций.
\end{enumerate}

\section{npm как система управления пакетами и сценариями сборки}

\subsection{Назначение npm}

\textbf{npm} --- пакетный менеджер для экосистемы Node.js. Он используется для:
\begin{itemize}
    \item установки зависимостей;
    \item публикации пакетов;
    \item запуска сценариев;
    \item управления версиями зависимостей.
\end{itemize}

Хотя npm не является системой сборки в строгом смысле, он часто управляет запуском сборочных инструментов: TypeScript, webpack, Vite, Rollup, ESLint и других.

\subsection{package.json}

Файл \texttt{package.json} описывает проект.

Пример:

\begin{lstlisting}
{
  "name": "demo",
  "version": "1.0.0",
  "scripts": {
    "build": "tsc",
    "test": "jest",
    "start": "node dist/index.js"
  },
  "dependencies": {
    "express": "^4.18.0"
  },
  "devDependencies": {
    "typescript": "^5.0.0",
    "jest": "^29.0.0"
  }
}
\end{lstlisting}

\subsection{Сценарии npm}

Сценарии запускаются командой:

\begin{lstlisting}
npm run build
\end{lstlisting}

Если в \texttt{package.json} определено:

\begin{lstlisting}
"scripts": {
  "build": "tsc"
}
\end{lstlisting}

то команда \texttt{npm run build} выполнит \texttt{tsc}.

\subsection{package-lock.json}

Файл \texttt{package-lock.json} фиксирует конкретные версии установленных зависимостей и помогает воспроизводить установку.

\subsection{Преимущества npm}

\begin{enumerate}
    \item огромная экосистема пакетов;
    \item простое описание зависимостей;
    \item удобные сценарии;
    \item lock-файл;
    \item интеграция с Node.js.
\end{enumerate}

\subsection{Недостатки npm}

\begin{enumerate}
    \item большое число транзитивных зависимостей;
    \item возможны проблемы безопасности цепочек поставки;
    \item сборка часто зависит от внешних инструментов;
    \item семантические диапазоны версий могут приводить к неожиданным обновлениям.
\end{enumerate}

\section{Сравнение универсальных и языко-ориентированных систем}

\subsection{Универсальные системы}

К универсальным системам относятся:
\begin{itemize}
    \item Make;
    \item Ninja;
    \item SCons;
    \item Bazel;
    \item Buck.
\end{itemize}

Их особенности:
\begin{enumerate}
    \item не привязаны к одному языку;
    \item описывают общие зависимости и команды;
    \item могут использоваться для произвольных задач;
    \item требуют явного описания многих деталей.
\end{enumerate}

\subsection{Языко-ориентированные системы}

К языко-ориентированным относятся:
\begin{itemize}
    \item Cargo для Rust;
    \item Maven и Gradle для Java/JVM;
    \item npm/yarn/pnpm для JavaScript;
    \item pip/poetry для Python;
    \item Go modules для Go.
\end{itemize}

Их особенности:
\begin{enumerate}
    \item знают стандартную структуру проекта;
    \item интегрированы с компилятором или интерпретатором;
    \item управляют зависимостями;
    \item поддерживают публикацию пакетов;
    \item обычно проще для типовых проектов на своём языке.
\end{enumerate}

\subsection{Сравнительная таблица}

\begin{center}
\begin{tabular}{|p{3cm}|p{4cm}|p{4cm}|p{4cm}|}
\hline
Система & Основная область & Управление зависимостями & Особенности \\
\hline
Make & универсальная сборка & почти отсутствует & простая модель целей и команд \\
\hline
CMake & C/C++, кроссплатформенная сборка & частично, через find\_package и FetchContent & генератор систем сборки \\
\hline
Ninja & быстрый backend сборки & отсутствует & обычно генерируется другими инструментами \\
\hline
Cargo & Rust & встроенное & единый стандарт для Rust-проектов \\
\hline
Maven & Java/JVM & встроенное, транзитивное & POM, жизненный цикл, репозитории \\
\hline
Gradle & JVM, Android, многомодульные проекты & встроенное & гибкие скрипты, плагины, кэширование \\
\hline
npm & JavaScript/Node.js & встроенное & package.json, scripts, lock-файл \\
\hline
\end{tabular}
\end{center}

\section{Управление модулями и зависимостями}

\subsection{Основные задачи менеджера зависимостей}

Менеджер зависимостей решает следующие задачи:
\begin{enumerate}
    \item загрузка внешних библиотек;
    \item разрешение транзитивных зависимостей;
    \item выбор совместимых версий;
    \item кэширование пакетов;
    \item проверка целостности;
    \item публикация пакетов;
    \item поддержка lock-файлов.
\end{enumerate}

\subsection{Семантическое версионирование}

Часто используется схема:
\[
MAJOR.MINOR.PATCH.
\]

Например:
\[
2.5.1.
\]

Смысл:
\begin{itemize}
    \item \texttt{MAJOR} меняется при несовместимых изменениях API;
    \item \texttt{MINOR} меняется при добавлении совместимой функциональности;
    \item \texttt{PATCH} меняется при исправлении ошибок без изменения API.
\end{itemize}

\subsection{Диапазоны версий}

Менеджеры пакетов часто позволяют указывать не точную версию, а диапазон.

Примеры:
\begin{itemize}
    \item \texttt{1.2.3} --- точная версия;
    \item \texttt{\^{}1.2.3} --- совместимые обновления в пределах major-версии;
    \item \texttt{\~{}1.2.3} --- обычно обновления patch-уровня;
    \item \texttt{>=1.2.0 <2.0.0} --- явный диапазон.
\end{itemize}

\subsection{Lock-файлы}

\textbf{Lock-файл} --- файл, фиксирующий точные версии всех зависимостей, включая транзитивные.

Примеры:
\begin{itemize}
    \item \texttt{Cargo.lock};
    \item \texttt{package-lock.json};
    \item \texttt{yarn.lock};
    \item \texttt{pnpm-lock.yaml};
    \item \texttt{gradle.lockfile}.
\end{itemize}

Lock-файл повышает воспроизводимость сборки.

\section{Алгоритм разрешения зависимостей}

\subsection{Общая постановка}

Пусть имеется множество пакетов. Каждый пакет может иметь несколько версий. Каждая версия задаёт ограничения на версии других пакетов.

Требуется выбрать по одной версии каждого необходимого пакета так, чтобы все ограничения были выполнены.

Формально:
\[
P = \{p_1, p_2, \ldots, p_n\}
\]
--- множество пакетов.

Для каждого пакета $p_i$ есть множество доступных версий:
\[
V_i = \{v_{i1}, v_{i2}, \ldots, v_{ik}\}.
\]

Нужно найти отображение:
\[
f: P \rightarrow \bigcup_i V_i,
\]
такое что:
\[
f(p_i) \in V_i
\]
и все ограничения зависимостей выполнены.

\subsection{Упрощённый алгоритм разрешения зависимостей с backtracking}

\textbf{Шаги:}
\begin{enumerate}
    \item Начать с корневых зависимостей проекта.
    \item Выбрать пакет, для которого ещё не выбрана версия.
    \item Найти все версии, удовлетворяющие текущим ограничениям.
    \item Выбрать одну версию.
    \item Добавить её зависимости в множество ограничений.
    \item Если возник конфликт, откатиться к предыдущему выбору.
    \item Повторять, пока все версии не будут выбраны.
\end{enumerate}

\subsection{Псевдокод}

\begin{lstlisting}
resolve(requirements):
    solution = {}

    if search(requirements, solution):
        return solution
    else:
        error "no compatible versions"

search(requirements, solution):
    if all requirements satisfied:
        return true

    pkg = choose_unresolved_package(requirements, solution)

    for version in candidate_versions(pkg, requirements):
        solution[pkg] = version
        new_requirements = requirements + dependencies(pkg, version)

        if consistent(new_requirements, solution):
            if search(new_requirements, solution):
                return true

        remove solution[pkg]

    return false
\end{lstlisting}

\subsection{Иллюстрация}

Пусть проект зависит от:
\[
A: [1.0, 2.0), \quad B: [1.0, 2.0).
\]

Доступные версии:
\[
A = \{1.0, 1.1\}, \quad B = \{1.0, 1.1\}.
\]

Дополнительно:
\[
A:1.1 \rightarrow C: [2.0,3.0),
\]
\[
B:1.1 \rightarrow C: [1.0,2.0).
\]

Если выбрать:
\[
A=1.1,\quad B=1.1,
\]
то возникает конфликт:
\[
C \in [2.0,3.0) \cap [1.0,2.0) = \varnothing.
\]

Алгоритм должен откатиться и попробовать другую версию, например:
\[
A=1.0,\quad B=1.1.
\]

\section{Теоретические замечания}

\subsection{Теорема о топологической сортировке}

\textbf{Теорема.}
Ориентированный граф имеет топологическую сортировку тогда и только тогда, когда он является ациклическим.

\subsubsection{Доказательство}

Докажем необходимость.

Пусть ориентированный граф $G$ имеет топологическую сортировку. Предположим противное: в $G$ есть ориентированный цикл:
\[
v_1 \rightarrow v_2 \rightarrow \cdots \rightarrow v_k \rightarrow v_1.
\]

По определению топологической сортировки для каждого ребра вершина-источник должна стоять раньше вершины, в которую ведёт ребро. Тогда:
\[
v_1 \text{ стоит раньше } v_2,
\]
\[
v_2 \text{ стоит раньше } v_3,
\]
\[
\cdots
\]
\[
v_k \text{ стоит раньше } v_1.
\]

Из первых $k-1$ неравенств следует, что:
\[
v_1 \text{ стоит раньше } v_k.
\]

Но из ребра
\[
v_k \rightarrow v_1
\]
следует, что:
\[
v_k \text{ стоит раньше } v_1.
\]

Получаем противоречие. Следовательно, циклов быть не может.

Докажем достаточность.

Пусть $G$ --- ориентированный ациклический граф. Покажем, что в нём существует вершина с нулевой входящей степенью.

Предположим противное: у каждой вершины есть хотя бы одно входящее ребро. Возьмём произвольную вершину $v_0$. Так как у неё есть входящее ребро, существует вершина $v_1$ такая, что:
\[
v_1 \rightarrow v_0.
\]
У вершины $v_1$ тоже есть входящее ребро, значит существует $v_2$:
\[
v_2 \rightarrow v_1.
\]

Продолжая этот процесс, получим бесконечную последовательность вершин:
\[
v_0, v_1, v_2, \ldots
\]

Но граф конечен, значит некоторая вершина повторится. Это образует ориентированный цикл, что противоречит ацикличности. Следовательно, существует вершина с нулевой входящей степенью.

Удалим эту вершину из графа вместе с исходящими рёбрами. Оставшийся граф также ацикличен. Повторяя рассуждение, можно последовательно удалять вершины с нулевой входящей степенью.

Порядок удаления вершин образует топологическую сортировку.

Теорема доказана.

\subsection{Следствие для систем сборки}

Если граф зависимостей целей ацикличен, то существует корректный порядок сборки. Если в графе есть цикл, то корректный порядок сборки невозможен без дополнительных соглашений или изменения структуры зависимостей.

\section{Практические рекомендации}

\subsection{Для небольших C/C++ проектов}

Можно использовать Make, если:
\begin{itemize}
    \item проект небольшой;
    \item целевая платформа Unix-подобная;
    \item зависимости просты;
    \item требуется минимальная инфраструктура.
\end{itemize}

\subsection{Для переносимых C/C++ проектов}

Лучше использовать CMake, если:
\begin{itemize}
    \item проект должен собираться на Linux, macOS и Windows;
    \item нужны разные генераторы;
    \item есть библиотеки и исполняемые файлы;
    \item требуется установка, тестирование, экспорт целей.
\end{itemize}

\subsection{Для Rust}

Обычно следует использовать Cargo, поскольку он является стандартным инструментом экосистемы Rust.

\subsection{Для Java}

Для классических Java-проектов часто используют Maven. Для Android, Kotlin и сложных многомодульных проектов часто выбирают Gradle.

\subsection{Для JavaScript и TypeScript}

npm, yarn или pnpm используются для управления пакетами, а непосредственную сборку выполняют специализированные инструменты: TypeScript compiler, webpack, Vite, Rollup и другие.

\section{Заключение}

Системы сборки и управления модулями решают фундаментальную задачу автоматизации получения артефактов из исходных файлов и зависимостей. Их основой является граф зависимостей, по которому определяется порядок сборки и необходимость пересборки.

Универсальные системы, такие как Make и Ninja, предоставляют низкоуровневую модель целей и команд. CMake занимает промежуточное положение: он описывает проект на более высоком уровне и генерирует файлы для конкретной системы сборки. Языко-ориентированные инструменты, такие как Cargo, Maven, Gradle и npm, тесно интегрированы со своими экосистемами и дополнительно решают задачу управления пакетами и версиями.

Выбор инструмента зависит от языка, размера проекта, требований к переносимости, воспроизводимости, скорости сборки и удобства сопровождения.

\end{document}
